\subsection{Smith Normal Form and Module Presentations}

\begin{perspective} \label{per:smith-normal-form-and-module-presentations} \hlt{Smith Form as Diagonalized Relations}\\
For a presentation $R^n\xrightarrow{A}R^m\to M\to0$, Smith form replaces $A$ by a diagonal relation matrix, decoupling finite relations into cyclic constraints.
\end{perspective}

\begin{theorem} \label{thm:smith-normal-form-over-pid} \hlt{Smith Normal Form over a PID}\\
Let $A\in\operatorname{Mat}(m,n,R)$, viewed as $R^n\to R^m$.\\
$\exists$ invertible $P\in\operatorname{Mat}(m,m,R)$ and $Q\in\operatorname{Mat}(n,n,R)$ such that $PAQ=\operatorname{diag}(d_1,\ldots,d_r,0,\ldots,0)$,\\
where $d_i\in R\setminus\{0\}$, $d_1\mid\cdots\mid d_r$, and $r=\operatorname{rk}_R(\operatorname{im}A)$. The ideals $\langle d_i\rangle$ are unique.
\begin{proof}
\begin{enumerate}[label=(\roman*)]
\setlength{\itemsep}{0pt}
\item \hlt{Existence}: If $A=0$, done. Set $I(A)=\langle a_{ij}\rangle=\langle d\rangle$. Elementary operations move a Bezout combination equal to $d$ to position $(1,1)$. If a pivot $p\nmid a$, Bezout on $(p,a)$ replaces $p$ by $\gcd(p,a)$, so $\langle p\rangle$ strictly increases; ACC in a PID forces termination. Then $p$ divides every entry, so clear its row and column.\\
Induct on the complementary submatrix. Finally replace each bad adjacent pair $(a,b)$ by $(\gcd(a,b),\operatorname{lcm}(a,b))$ until $d_1\mid\cdots\mid d_r$.
\item \hlt{Uniqueness}: Let $\Delta_k(A)$ be the ideal of $k\times k$ minors and $\Delta_0(A)=R$. Matrix equivalence preserves each $\Delta_k$. In Smith form, $\Delta_k(A)=\langle d_1\cdots d_k\rangle$ for $1\leq k\leq r$, and $\Delta_k(A)=0$ for $k>r$. Since $R$ is a domain, $(\Delta_k(A):\Delta_{k-1}(A))=\langle d_k\rangle$, so $d_k$ is determined up to associates.
\end{enumerate}
\noindent
\hlt{Structural Commentary}: Smith form preserves rank plus divisibility data.\\
\hlt{Key Insight}: Determinantal ideals encode invariant factors.
\end{proof}
\end{theorem}

\begin{theorem} \label{thm:adapted-basis-submodule-free-pid} \hlt{Adapted Bases for Submodules of Free Modules}\\
Let $N\leq F$, $\operatorname{rk}_R(F)=n$, and $\operatorname{rk}_R(N)=k$.\\
Then $\exists(v_i)_{i=1}^n\in\operatorname{Bas}_R(F)$ and $r_1,\ldots,r_k\in R\setminus\{0\}$ such that $r_1\mid\cdots\mid r_k$ and $(r_iv_i)_{i=1}^k\in\operatorname{Bas}_R(N)$.
\begin{proof}
Choose bases of $N$ and $F$, and let $A$ represent $\iota:N\hookrightarrow F$.\\
Smith form gives $PAQ=\operatorname{diag}(r_1,\ldots,r_k,0,\ldots,0)$ with $r_i\neq0$ and $r_1\mid\cdots\mid r_k$, since $\iota$ is injective.\\
The row operations give a basis $(v_i)_{i=1}^n$ of $F$ and the column operations give a basis $(w_i)_{i=1}^k$ of $N$ with $\iota(w_i)=r_iv_i$. Hence $(r_iv_i)_{i=1}^k=(w_i)_{i=1}^k$ is a basis of $N$.\\
\hlt{Structural Commentary}: Inclusions diagonalize after choosing bases.\\
\hlt{Key Insight}: Diagonal entries are the required rescalings.
\end{proof}
\end{theorem}

\begin{corollary} \label{cor:smith-normal-form-presents-pid-modules} \hlt{Smith Form Presents Finitely Generated PID Modules}\\
Let $R^n\xrightarrow{A}R^m\to M\to0$ be a finite presentation.\\
If $A\sim\operatorname{diag}(d_1,\ldots,d_r,0,\ldots,0)$ in Smith form. $I=\{i:1\leq i\leq r,\ d_i\notin R^\times\}$, then $M\cong R^{m-r}\oplus\bigoplus_{i\in I}R/\langle d_i\rangle$.
\begin{proof}
Equivalence $A\sim PAQ$ changes bases in $R^m,R^n$, so $\operatorname{coker}A\cong\operatorname{coker}(PAQ)$.\\
For the diagonal map, $d_i\in R^\times$ gives $0$, $d_i\notin R^\times$ gives $R/\langle d_i\rangle$, and each zero row gives $R$.\\
\hlt{Structural Commentary}: Smith form turns presentations into direct sums.\\
\hlt{Key Insight}: Diagonal relations create cyclic torsion summands, and zero rows create free summands.
\end{proof}
\end{corollary}

\begin{example} \label{ex:smith-normal-form-abelian-group-decomposition} \hlt{Smith Form for a Finite Abelian Group}\\
Let $M=\mathbb{Z}^2/\langle\!\langle(6,0),(0,10),(2,4)\rangle\!\rangle_{\mathbb{Z}}$, with relation matrix $A=\begin{bmatrix}6&0&2\\0&10&4\end{bmatrix}$. Then $\Delta_1(A)=\langle6,10,2,4\rangle=\langle2\rangle$ and $\Delta_2(A)=\langle60,24,-20\rangle=\langle4\rangle$. Hence $d_1\sim2$, $d_1d_2\sim4$, so $d_2\sim2$ and $M\cong\mathbb{Z}/2\mathbb{Z}\oplus\mathbb{Z}/2\mathbb{Z}$.\\
\hlt{Key Insight}: Determinantal ideals read off cyclic summands.
\end{example}
